\section{Mathematical results preserved from reports 64--69}
\label{sec:completion-theorems}

The six deliveries in arrival commit \code{fe88ae28e} contain mathematical
results beyond implementation proposals. Their complete manuscripts, PDFs,
proof appendices, source code and experimental records are preserved in
\code{reports/64} through \code{reports/69}. This section records their
mathematical content independently of whether a proposed implementation is
promoted. A slower interface, a conditional recognition theorem, or a
counterexample remains worth preserving. The next section treats the one
runtime change selected from this batch.

All six ZIPs passed CRC validation. Placement removed only the common archive
wrapper and redundant pure checksum manifests; text line endings were
normalized. The retained placement record hashes every preserved file.
Reports 64 and 69 also supplied whole-package SHA256 manifests, checked before
normalization. The original arrival commit retains the original archives.
There are 33, 33, 45, 649, 652 and 702 preserved files respectively, including
historical source snapshots. These snapshots are evidence, not replacements
for the maintained runtime. The report catalogue links each master manuscript.

\subsection{A surface gluing identity and a disc criterion}

Reports 64 and 65 work with $r\ge1$ individually labelled, collared boundary
intervals. Distinct intervals are disjoint including endpoints. Every
fragment component meets an interval, has nonempty boundary, and retains
boundary arcs outside the seam intervals. The quotient must be a surface.
Let $P,Q$ partition the interval labels according to their component in the
two fragments $S,T$. Form a bipartite multigraph $H$: its vertices are the
blocks of $P$ and $Q$, and each label gives its own seam edge. Parallel edges
must not be suppressed.

\begin{proposition}[Boundary-interval gluing]
Under these hypotheses, the union $U$ satisfies
\begin{align*}
 c(U)&=|P\vee Q|, & \chi(U)&=\chi(S)+\chi(T)-r,\\
 \beta_1(U;\mathbb F_2)&=\beta_1(S;\mathbb F_2)+\beta_1(T;\mathbb F_2)
       +r-|P|-|Q|+|P\vee Q|.
\end{align*}
In particular $U$ is a disc exactly when all fragment components are discs
and $H$ is a tree.
\end{proposition}
\begin{proof}
Connected components are exactly the connected components of $H$, giving
the join formula. Each interval has Euler characteristic one, so cellular
inclusion--exclusion gives the Euler formula. A compact surface all of whose
components have boundary has $H_2(-;\mathbb F_2)=0$ and hence
$\beta_1=c-\chi$. Substitution proves the third identity; its last term is
the cycle rank of $H$ and is nonnegative. If $U$ is a disc, all three
nonnegative contributions vanish and $H$ is connected. A connected compact
surface with boundary and first Betti number zero is a disc, so the fragment
components are discs and $H$ is a tree. Conversely remove leaves of that
tree: each step glues a disc along one boundary interval to another disc.
The result remains a disc, irrespective of the interval parametrization.
\end{proof}

Thus the defect $\delta=c-\chi$ is additive with the graph cycle correction.
Positive defect can be pruned only while the continuation obeys this seam
contract. Later capping or cutting changes the problem: capping one boundary
of an annulus gives a disc. Cyclic order and twists disappear from the tree
criterion because it forbids cycles; they do not disappear from arbitrary
surface gluing. A binary run describing $2^B$ parallel seams is not one of
these individually labelled intervals. The width parameter is the actual
interface, not the number of run records.

\subsection{Exact exterior completion ranks and original witnesses}

Set $n=r-1$ and delete the root coordinate from the incidence representation
of the complete graph on the labels, over $\mathbb F_2$. For a partition $P$,
choose a spanning tree in each block and wedge its incidence columns to form
$f_P\in\Lambda^{d(P)}(\mathbb F_2^n)$, where $d(P)=r-|P|$.
The construction is independent of the chosen block trees: all give the
same minors in characteristic two. A union of incidence columns is
independent exactly when its multigraph is a forest. Dependence gives an
even-degree subgraph containing a cycle; leaf removal proves the converse.
The graph obtained from the two block forests and the seam graph have the
same component and cycle counts. Consequently their union is a spanning
tree exactly when the gluing graph is a tree. Equivalently, for
$d(P)+d(Q)=n$, the top coefficient of $f_P\wedge f_Q$ is one exactly for a
disc completion.

\begin{proposition}[Sharp completion dimensions]
The disc-completion matrix has rank $2^{r-1}$ over $\mathbb F_2$.
Its grade-$d$ block has rank $\binom{r-1}{d}$.
\end{proposition}
\begin{proof}
The exterior factorization gives the upper bounds. For each subset
$A\subseteq\{1,\ldots,n\}$, take the partition with block $\{0\}\cup A$
and all other labels singleton. Its feature is the standard wedge $e_A$.
The complementary rooted-star partition has feature $e_{A^c}$; their
pairing is nonzero, and pairings with other subsets of the same grade vanish.
Thus the complementary-grade pairing is nondegenerate on the full standard
basis. Summing the grade dimensions gives $2^n$.
\end{proof}

For a weighted family, process candidates by increasing cost and retain a
candidate whenever its feature is independent of earlier retained features.
Every dropped feature is a linear combination of features of no greater
cost. Pairing with any completing feature gives one for a feasible dropped
candidate, so at least one summand is feasible and costs no more. Since the
retained family is a subfamily, it preserves the exact minimum for every
admitted completion. It retains at most $2^{r-1}$ \emph{original} witnesses.
Linear combinations certify domination; they do not describe new embedded
surfaces. Hidden geometric conditions must be uniform within each reduction
type or proved to respect the feature map.

Rank-based weighted connectivity and matroid representative families are
classical. Report 64 explicitly compares its construction with
Bodlaender--Cygan--Kratsch--Nederlof and the acyclic representatives of
Bergougnoux--Kant\'e, Theorem 3.8
\cite{completion-rank-classical,completion-acyclic-classical}.
The sharp dimensions and their surface contract are the identities reviewed
here; this is not a claim that rank-based compression itself is new.

\subsection{Why full cut strata are larger}

Report 64 also proves an exact comparison that need not be used by the
runtime to remain mathematically useful. Anchor a cut variable $x_0=0$ and
work with Boolean square-free polynomials over $\mathbb F_2$. The cut indicator
of $P$ is
\[
 C_P(x)=\prod_{uv\text{ in a block forest of }P}(1+x_u+x_v).
\]
Its degree-$d(P)$ homogeneous part is $f_P$: choosing distinct variables
from all factors gives the incidence permanent, equal to the determinant
in characteristic two. Repeated variables contribute only lower degrees.

For $0\le d\le r-2$, the span $W_d$ of grade-$d$ cut functions is exactly
the space of square-free polynomials of degree at most $d$. To see the
inclusion $W_{d-1}\subseteq W_d$, a grade-$(d-1)$ partition has at least
three blocks. Merge three chosen block pairs $ab,ac,bc$. Their indicators
sum to the original indicator because the three equality factors sum to one.
The rooted-star function $\prod_{i\in A}(1+x_i)$ has leading term $x_A$.
Induction therefore supplies every monomial of degree at most $d$. The
reverse inclusion follows from the degree bound. At grade $r-1$ there is
just the one-block partition, with rank one. Hence
\[
 \operatorname{rank} C_{r,d}=
 \begin{cases}\sum_{j=0}^d\binom{r-1}{j},&0\le d\le r-2,\\
 1,&d=r-1.
 \end{cases}
\]
Grade-separated full cut bases retain $(r-1)2^{r-2}+1$ rows for $r\ge2$,
whereas the homogeneous exterior bases retain $2^{r-1}$. Both remain single
exponential. Projection removes information irrelevant to the complementary
grade rather than changing the asymptotic class. A Boolean subset M\"obius
transform and grade projection convert an existing cut table in
$O(r2^{r-1})$ binary scalar operations.

\subsection{Signed compatibility and the finite-group formula}

Report 66 adds orientation consistency to connectivity. A signed partition
imposes $x_i+x_j=\epsilon_i+\epsilon_j$ within each block. Anchor $x_0=0$
and use the indicator vector of its satisfying assignments. The inner
product of two such vectors counts their common assignments modulo two.
An inconsistent union has none; a consistent union with $c$ blocks has
$2^{c-1}$. This is one exactly when the union is connected. Thus the signed
compatibility matrix also has rank $2^{r-1}$: one-block signed states give
an identity submatrix. The same cost-ordered original-witness argument
preserves minimum cost, including negative costs. It does not preserve the
number of feasible witnesses.

The full finite-group result is stronger and is preserved even though no
native use is yet selected. Let $G$ have order $q$. A gain partition labels
each block by $\alpha_i\in G$ modulo simultaneous \emph{left} multiplication
and imposes $x_j=x_i\alpha_i^{-1}\alpha_j$. This convention works for
noncommutative groups. There is no zero block. The number of states is
\[
 D_r(q)=\sum_{k=1}^r S(r,k)q^{r-k},
\]
where $S(r,k)$ is a Stirling number of the second kind. Let $Z$ be the zeta
matrix for extension of these constraints. Define the compatibility kernel
$K_G$ to be one for consistent connected unions and zero otherwise.

\begin{proposition}[Integer factorization and field-dependent rank]
With $k(s)$ the number of blocks of state $s$,
\begin{align*}
 K_G&=Z\operatorname{diag}(g(s))Z^{\mathsf T},
 & g(s)&=(-q)^{k(s)-1}(k(s)-1)!,\\
 \det K_G&=\prod_{k=1}^r
       \bigl((-q)^{k-1}(k-1)!\bigr)^{S(r,k)q^{r-k}}.
\end{align*}
Its rank is $D_r(q)$ in characteristic zero. In characteristic $p>0$ it is
\[
 \begin{cases}
 q^{r-1},&p\mid q,\\
 \sum_{k=1}^{\min(r,p)}S(r,k)q^{r-k},&p\nmid q.
 \end{cases}
\]
\end{proposition}
\begin{proof}
Let $f$ indicate one-block states. A $k$-block state has $q^{k-1}$ one-block
extensions. For each fixed extension, its interval is the ordinary partition
lattice on those $k$ blocks: the final gauges fix all intermediate labels.
Its M\"obius value is $(-1)^{k-1}(k-1)!$. M\"obius inversion of $f$ thus
gives $g(s)$ above. Summing $g$ over common extensions of two states yields
one for their consistent connected union and zero for every other union,
giving the matrix identity. The zeta matrix is unitriangular over every
field. Rank is therefore the number of nonzero diagonal entries, and its
determinant is their product. In characteristic $p\mid q$ only $k=1$
survives; otherwise $(k-1)!$ survives exactly for $k\le p$.
\end{proof}

Already the $r=2,q=2$ determinant is $-2$: the rank is two over
$\mathbb F_2$ and three over $\mathbb Q$. For fixed $q$, the rational rank
is $2^{\Theta(r\log r)}$, precluding a uniformly single-exponential
\emph{rational bilinear representation of this kernel}. This does not rule
out other algorithms or more restricted continuation families. The
factorization is an incidence-algebra argument in the gain-partition setting;
the original manuscript records the classical Dowling and join-matrix context.

For actual orientable fragments, signed seam consistency encodes whether
their orientations extend. Euler characteristic adds with the known seam
correction. Cost $-\chi$ therefore preserves maximum Euler characteristic
among connected orientable completions. A permanent unglued boundary anchor
ensures nonempty boundary, where $\chi\le1$ with equality precisely for a
disc. The delivered complete \emph{supplied} triangulated-patch selection
algorithm costs $\operatorname{poly}(N,B,b)2^{3b}$ elementary bit operations
for frontier width $b$. Whole-site symbolic transfers avoid allocating the
temporarily larger exponential space. Reports 64 and 65 give corresponding
typed assembly contracts. These are meaningful finite-input algorithms.
General knot recognition still needs a source-complete geometric patch
producer, controlled width, embeddedness and essential-boundary interfaces.
The maintained diagram-to-exterior construction already exists; it does
not by itself supply that assembly theorem.

\subsection{Farkas certificates and a propagation obstruction}

Report 67 treats canonical standard coordinates by choosing a zero triangle
anchor at every global vertex. In a fixed allowed sector $J$, positive
canonical Euler characteristic is rational feasibility of
\[
 M_Jz=0,\qquad z\ge0,\qquad c_J^{\mathsf T}z\ge1.
\]
The scaling inequality is exact, not a numerical positivity threshold.
A negative homogeneous Farkas certificate is a rational $y$ with
$M_J^{\mathsf T}y\ge c_J$: every feasible nonnegative matching vector then
has $c_J^{\mathsf T}z\le0$. The dual alternative gives completeness for
this homogeneous linear problem. Anchor enumeration covers the canonical
cone. A componentwise objective envelope may prove a negative result, but
a positive envelope must be checked against an actual anchor objective.

If forbidding a type makes the positive polyhedron empty, that type is
mandatory. Every admissible positive witness must then set its two
incompatible peers to zero. Repeated deletions preserve every admissible
positive witness; exhaustive propagation has an order-independent terminal
closure. A strong mandatory-only backdoor requires this closure to leave
compatible allowed masks for every anchor and every assignment on its
chosen tetrahedra. This is stronger than finding an admissible basic vector.

For a canonical admissible positive-Euler vector $x$, let $Z(x)$ be its
tetrahedra with all three quadrilateral coordinates zero, and put
$U(T)=\bigcup_x Z(x)$. Report 67 proves
\[
 U(T)\subseteq B,\qquad
 \beta_{\rm LP}(T)\ge |U(T)|\ge t-|\operatorname{qsupp}(x)|
\]
for every strong mandatory-only backdoor $B$. Choose anchors and assignments
consistent with $x$. It survives every deletion. At a tetrahedron of $Z(x)$
outside $B$, $x$ witnesses feasibility with each type forbidden, so none
can become mandatory. No peer deletion occurs there and all three types
remain allowed, contradicting terminal mask compatibility. Taking the union
proves the obstruction. Sparse positive witnesses can thus force a large
strong parameter, even when they are immediately easy to find.

A policy-dependent deciding set permits early termination on an actual
admissible positive vector. It avoids that mask-only obstruction on positive
instances, but on instances with no canonical admissible positive-Euler
vector it equals the strong parameter. That negative condition is not the
same as knottedness. With a polynomial-time exact LP policy, a supplied
deciding set of size $b$ gives $A(T)3^b\operatorname{poly}(t)$ work; search
for an unknown set costs
$A(T)\sum_{j\le b}\binom tj3^j\operatorname{poly}(t)$.
For one vertex $A(T)=4t$. The delivered Bland simplex has no proved
polynomial pivot bound, and no general logarithmic deciding-set theorem
is supplied. The manuscript also proves that the whole Fibonacci layered
family has one positive canonical admissible meridian ray, full
quadrilateral support, and mandatory strong backdoor zero. Large support
and easy discovery are compatible; the family does not establish a
universal parameter bound. These proofs and negative certificates are
preserved without installing the larger LP overlay over the newer runtime.

\subsection{Exact transport changes and disconnected primitive fibres}

Report 68 analyzes a legal relative $2$--$3$ bipyramid move for a transported
integral cocycle. Sort the three belt heights as $x\le y\le z$ and the
two apex heights as $L\le U$. Set $[a]_+=\max(a,0)$ and
\begin{align*}
 g&=[L-z]_+ +[x-U]_+,\\
 \Delta P&=\max(U,y)-\min(L,y)
       +[\min(U,x)-L]_+ +[U-\max(L,z)]_+.
\end{align*}
Then the upward move changes Euler characteristic by $-2g$ and the number
of normal pieces by $\Delta P\ge U-L\ge0$. The downward move has the
opposite changes. The proof counts relative cells: the vertex change is
$|A-B|$, the internal arc change is
$\sum_{C,D,E}\operatorname{range}(A,B,C)-\operatorname{range}(C,D,E)$,
and the piece change is the difference of the five tetrahedral height
spans. Unchanged boundary cells cancel. Sorting the heights reduces these
span expressions to the displayed formulas. A new upward edge difference
is at most twice the previous maximum coefficient; downward removal adds
none. Thus coefficient-bit growth can be charged to upward moves.

Crucially, preserving a primitive cohomology class does not preserve
connectedness of its fibre. The retained source-bound example changes a
connected 24-piece essential disc of Euler characteristic one into an
11-piece disc and two four-piece spheres. The total has three components,
19 pieces and Euler characteristic five. Belt/apex heights are
$(0,2,1,5,4)$: $g=2$ and $\Delta P=5$ for the reversed upward move.
Fresh replay checks both source-bound moves, both archived component
certificates and both current topology-spectrum certificates. An independent
cell-span calculation checks all $5^5=3125$ small height assignments.
Primitive \emph{matching-coordinate rays}, used in the next section, have
a different connectedness proof. They must not be conflated with primitive
cohomology fibres.

There are two further useful results even without native promotion.
First, let $m_v$ be the minimum triangle coordinate at global vertex $v$.
During a fixed-state scoring pass, the thirteen least corner records per
vertex suffice to find each new minimum after a legal $3$--$2$ move:
at most twelve old corners are removed, so a surviving minimum occurs among
the first thirteen; compare it with the at most eight new corners. This
gives linear preparation and constant integer work per candidate. Summaries
must be rebuilt after committing a move. Global regauging is outside the
lemma's contract. The refined scorer is proved but unimplemented.

The same report proves $m'_v\ge m_v$ and nonincrease of the piece count
after vertex-link peeling. Write
$a=[\min(U,x)-L]_+$ and $b=[U-\max(L,z)]_+$.
The local apical minima rise by $g$ and belt minima remain unchanged.
Taking minima over identified vertices and unchanged exterior corners gives
$m'_v\ge m_v$. If $r_v$ is the number of formal apices identified with $v$,
then $|C'_v|=|C_v|-2r_v$, $\sum_v r_vm_v\le a+b$, and
$\Delta P\ge2(a+b)$. With $\delta_v=m'_v-m_v\ge0$ the peeled change is
\[
 -\Delta P+2\sum_v r_vm_v-\sum_v\delta_v|C'_v|\le0.
\]
This excludes gcd division and global regauging. It proves neither peeled
Euler monotonicity nor bounded essential-disc discovery.

Second, cutting a marked degree-two boundary curve leaves paths with exactly
two endpoint \emph{occurrences}. Four additive weights---count, sum of
endpoint labels, sum of squared labels, and length---recover a path's
endpoints: if $S=x+y$ and $Q=x^2+y^2$, then
$(x-y)^2=2Q-S^2$. Integral recovery and endpoint adjacency restore the cyclic
order. This is a degree-two interface theorem, not reconstruction of
arbitrary attachments from moments. The report's broader transport search
bounds require explicit coverage assumptions; its surveyed recognition
coverage does not increase. The complete theorem statements and negative
examples remain in the preserved manuscript.

\subsection{Support projections and the number of component types}

Report 69 fixes the positive support $S$ of a supplied admissible matching
vector and its rational matching nullity $d$. Projection to coordinates $J$
is injective on the supported kernel exactly when the complementary columns
are independent. A minimum-cost coordinate projection therefore retains
exactly $d$ coordinates: maximize the total weight of the complementary
column basis by the weighted matroid greedy rule. Here costs are allocated
binary slot sizes. A rational decoder and a checked nonsingular modular
minor certify exact rank and kernel reconstruction. This optimality is
among coordinate projections with a linear decoder on the entire supported
space; it is not an optimality claim among all compressed encodings.
Under uniform positive scaling, reusing the projection costs within $d$
bits of the new optimum because each slot gains the same scale bit length
up to one bit. Capacity-based scalar packing preserves equality of every
live point-weight run of this selected-vector observer, not merely its
final values. It does not prove fewer orbit runs. The full projection
interface has measured regressions and remains optional research code.

\begin{proposition}[Unmarked normal component-type inventory]
For the report's compact orientable manifold setting, let $H(x)$ count
distinct full normal-coordinate component types and let $a$ count supported
vertex-link types. Then $H(x)\le2d-a$. If all components are two-sided,
$H(x)\le d$.
\end{proposition}
\begin{proof}
Distinct mutually disjoint two-sided normal types have independent rational
coordinate vectors. A dependence can be separated into two positive integer
combinations. Both are realized by disjoint parallel copies of the selected
two-sided surfaces. Equal coordinates imply normal isotopy and hence equal
component-type multisets, contradicting their disjoint type sets.

For each nonlink type choose one component. Replace each one-sided component
by its connected two-sided horizontal interval-neighbourhood frontier,
whose vector is twice the original. Choose neighbourhoods disjointly and
coalesce equal resulting vectors, giving $h$ nonlink two-sided types. Each
has at most two original preimages, $z$ and $z/2$. It cannot become a
vertex link because its positive quadrilateral survives doubling. Along
with the $a$ link vectors, independence gives $h+a\le d$, while
$H-a\le2h$. Combining proves $H\le2d-a$.
\end{proof}

The bound counts coordinate types, not components or marked attachments.
The one-tetrahedron M\"obius source $3M$ has $d=1,a=0,H=2$; adding its
disjoint vertex link gives $d=2,a=1,H=3$. Both attain the bound. Binary
multiplicities remain essential. No global hierarchy branching bound follows
without a theorem controlling the markings.

The report also gives safe geometric support blocks: connect positive types
using their occurrences in each \emph{original, uncancelled} face equation.
Restricting the source to a block preserves matching, and actual face
attachments never leave a block. This gives a disjoint surface decomposition
and a direct sum of matching spaces. A graph built only after algebraic
cancellation can omit geometric attachments and needs a separate proof.
Its unit-pivot specialization is the integration in the next section.

\subsection{Fresh validation and limits of the review}

Scratch copies of reports 64--66 pass 18, 36 and 55 test methods respectively.
Their fresh audits reproduce: report 64's 44,168 partition pairs, 1,155
incidence minors, 3,336 weighted queries, 72 independent certificates,
2,000 ribbon surfaces, 5,295 projections and 1,800 staged joins; report
65's 813,297 direct matrix entries, 14,014 completions from 1,000 weighted
families, 1,000 literal triangulated gluings and 200 grammar comparisons;
report 66's 68,581 signed entries, 442 triangulated gluings and all 297
assignments in the finite patch system. The transport audit is described
above. These finite checks supplement the manuscript proofs; they cannot
establish their infinite statements by exhaustion. No fresh execution of
the complete report-67 LP overlay is claimed. Its original evidence is
preserved and its polynomial-policy hypothesis remains explicit.

The mathematical review accepts the statements under their stated contracts,
without claiming formal proof-assistant verification or originality beyond
the manuscripts' attributions. All original proofs remain available even
where this synthesis gives only their central argument. The generic LP,
geometric transport, signed-patch and packed-coordinate overlays have not
been installed over the newer native tree. Neither these mathematical
results nor the local speedup below supplies a general quasipolynomial or
general subexponential unknot-recognition bound.
